% !TEX root = ../main.tex
%% Chapitre 17 — Loi binomiale
%% Sources : frontend/cours/resume/terminal/resume_loi_binomiale_terminal.txt
%%           frontend/exercices/terminal_spé_Maths/exercice_loi_binomiale.txt
%% Exercice 3, question 5 : remplacée (elle reprenait la question 2).
%% Exercices 12, 13 et 14 : problèmes inédits rédigés pour le livre.

\begin{ouverture}[binomiale]{17}{Loi binomiale}{Partie 3 · Probabilités}
  {Épreuve de Bernoulli\enspace·\enspace Schéma de Bernoulli\enspace·\enspace Coefficients binomiaux\enspace·\enspace
   Espérance\enspace·\enspace Fluctuation}
  \objectif{Reconnaître une situation modélisée par une loi binomiale et la justifier.}
  \objectif{Calculer $P(X=k)$, $P(X\leq k)$, $P(X\geq k)$ et $P(a\leq X\leq b)$.}
  \objectif{Calculer et interpréter l'espérance, la variance et l'écart-type d'une loi binomiale.}
  \objectif{Déterminer un seuil ou un intervalle à l'aide de la calculatrice ou d'un programme Python.}
  \objectif{Utiliser un intervalle de fluctuation pour prendre une décision.}
  \prerequis{Variables aléatoires, espérance et variance (Première).}
  \prerequis{Coefficients binomiaux (chapitre 15) ; indépendance (chapitre 16).}
\end{ouverture}

%% =====================================================================
%%  L'ESSENTIEL
%% =====================================================================
\essentiel

\essSection{1}{Épreuve et loi de Bernoulli}

\begin{definition}[Épreuve de Bernoulli]
Une \textbf{épreuve de Bernoulli} est une expérience aléatoire à deux issues : le \emph{succès} $S$, de probabilité $p$, et
l'\emph{échec} $\overline S$, de probabilité $1-p$. La variable $X$ qui vaut $1$ en cas de succès et $0$ sinon suit la
\textbf{loi de Bernoulli} de paramètre $p$ :\qquad $E(X)=p$\qquad et\qquad $V(X)=p(1-p)$.
\end{definition}

\essSection{2}{Loi binomiale}

\begin{definition}[Schéma de Bernoulli et loi binomiale]
Un \textbf{schéma de Bernoulli} de paramètres $n$ et $p$ est la répétition de $n$ épreuves de Bernoulli \textbf{identiques et
indépendantes} de paramètre $p$. La variable $X$ égale au nombre de succès suit la \textbf{loi binomiale} $\mathcal{B}(n\pv p)$.
\end{definition}

\noindent
\begin{minipage}[c]{0.56\linewidth}
\begin{theoreme}[Probabilités]
\raggedright
Si $X$ suit la loi $\mathcal{B}(n\pv p)$, pour tout entier $k$ compris entre $0$ et $n$ :
\[ P(X=k)=\binom nk\,p^{k}\,(1-p)^{n-k}. \]
\end{theoreme}
\begin{demonstration}
\raggedright
Un chemin de l'arbre qui comporte $k$ succès et $n-k$ échecs a pour probabilité $p^{k}(1-p)^{n-k}$, par indépendance. Un tel
chemin est déterminé par le choix des $k$ épreuves réussies parmi les $n$ : il y en a $\binom nk$.
\end{demonstration}
\end{minipage}\hfill
\begin{minipage}[c]{0.4\linewidth}
\centering
\begin{tikzpicture}[otArbre, level distance=15mm, level 1/.style={sibling distance=20mm},
  level 2/.style={sibling distance=10mm}, level 3/.style={sibling distance=5mm},
  every node/.style={font=\scriptsize}]
  \node {}
    child {node {$S$}
      child {node {$S$} child {node (a) {$S$}} child {node (b) {$\overline S$}}}
      child {node {$\overline S$} child {node (c) {$S$}} child {node (d) {$\overline S$}}}}
    child {node {$\overline S$}
      child {node {$S$} child {node (e) {$S$}} child {node (f) {$\overline S$}}}
      child {node {$\overline S$} child {node (g) {$S$}} child {node (h) {$\overline S$}}}};
  \foreach \nd/\kk in {a/3,b/2,c/2,d/1,e/2,f/1,g/1,h/0}{%
    \node[otFeuille] at ([xshift=1mm]\nd.east) {$k=\kk$};}
\end{tikzpicture}\\[0.5mm]
\legende{$n=3$ : $\binom32=3$ chemins comportent $2$ succès.}
\end{minipage}

\essSection{3}{Espérance, variance et allure de la loi}

\noindent
\begin{minipage}[c]{0.46\linewidth}
\begin{theoreme}[Espérance et variance]
\raggedright
Si $X$ suit la loi $\mathcal{B}(n\pv p)$ :\\[1.5mm]
{\centering $E(X)=np$\qquad $V(X)=np(1-p)$\\[1mm] $\sigma(X)=\sqrt{np(1-p)}$\par}\vspace{1.5mm}
\emph{Démonstration au chapitre 18 : $X$ est la somme de $n$ variables de Bernoulli indépendantes.}
\end{theoreme}
\end{minipage}\hfill
\begin{minipage}[c]{0.5\linewidth}
\centering
\begin{tikzpicture}
\begin{axis}[otSchema, width=76mm, height=44mm, ybar, bar width=1.7mm, xmin=-0.8, xmax=15.8, ymin=0, ymax=0.22,
  xtick={0,2,...,14}, ytick={0.1,0.2}, yticklabels={$0{,}1$,$0{,}2$}, xlabel={$k$},
  xlabel style={at={(1,0)},anchor=west}, grid=major, grid style={otRule!70,line width=0.3pt}]
  \addplot[fill=otNavySoft,draw=otNavy,line width=0.4pt] coordinates {(0,0.0008) (1,0.0068) (2,0.0278) (3,0.0716) (4,0.1304)
    (5,0.1789) (6,0.1916) (7,0.1643) (8,0.1144) (9,0.0654) (10,0.0308) (11,0.0120) (12,0.0039) (13,0.0010) (14,0.0002) (15,0)};
  \draw[otGreen,dashed,line width=0.7pt] (axis cs:6,0) -- (axis cs:6,0.215);
\end{axis}
\end{tikzpicture}\\[0.5mm]
\legende{Loi $\mathcal{B}(20\pv0{,}3)$ : les valeurs se concentrent autour de $E(X)=6$.}
\end{minipage}

\essSection{4}{Calculer avec la loi binomiale}

\begin{savoirfaire}[Justifier une loi binomiale]
On écrit : « On répète $n$ fois, de façon identique et indépendante, une épreuve à deux issues ; le succès « … » a pour
probabilité $p$. La variable $X$ compte le nombre de succès, donc $X$ suit la loi $\mathcal{B}(n\pv p)$. » Un tirage sans remise
dans une \textbf{grande} population est assimilé à un tirage avec remise.
\end{savoirfaire}

\begin{savoirfaire}[Probabilités cumulées]
La calculatrice donne $P(X=k)$ et $P(X\leq k)$. On en déduit $P(X\geq k)=1-P(X\leq k-1)$ et
$P(a\leq X\leq b)=P(X\leq b)-P(X\leq a-1)$. Pour un seuil du type « au moins un succès », on résout $1-(1-p)^{n}\geq c$ à l'aide du
logarithme. En Python, les fonctions ci-dessous calculent $P(X=k)$ et $P(X\leq k)$ :
\end{savoirfaire}
\begin{python}
from math import comb

def P(n, p, k):     # probabilité P(X = k)
    return comb(n, k) * p**k * (1-p)**(n-k)

def F(n, p, k):     # probabilité cumulée P(X <= k)
    return sum(P(n, p, i) for i in range(k+1))
\end{python}

\begin{savoirfaire}[Intervalle de fluctuation]
Pour $X$ de loi $\mathcal{B}(n\pv p)$, on cherche le plus petit entier $a$ tel que $P(X\leq a)>0{,}025$ et le plus petit entier $b$
tel que $P(X\leq b)\geq0{,}975$ : alors $P(a\leq X\leq b)\geq0{,}95$. Si la valeur observée n'est pas dans $\intFF ab$, on rejette
l'hypothèse « $p$ est la bonne probabilité », avec un risque d'erreur inférieur à 5\,\%.
\end{savoirfaire}

\begin{attention}[Pièges fréquents]
\begin{itemize}[label={\color{otOrange}\textbullet},leftmargin=1.4em,itemsep=0.6mm,topsep=0.5mm]
  \item « Au moins $k$ » : $X\geq k$ ; « au plus $k$ » : $X\leq k$ ; « plus de $k$ » : $X>k$, soit $X\geq k+1$.
  \item $P(X\geq k)=1-P(X\leq k-1)$, et non $1-P(X\leq k)$.
  \item Sans indépendance des épreuves (tirages sans remise dans une petite urne), la loi n'est pas binomiale.
  \item L'espérance $np$ n'est pas forcément un entier : c'est une moyenne sur un grand nombre de répétitions.
\end{itemize}
\end{attention}

%% =====================================================================
%%  EXERCICES
%% =====================================================================
\exercices
\rubrique{Découverte}

\begin{exercice}[1]{Un jeu de cartes}
Un jeu de société contient $20$ cartes : $12$ bleues et $8$ rouges. On tire successivement $5$ cartes avec remise et l'on note $X$
le nombre de cartes bleues obtenues.
\begin{questions}
  \item Justifier que $X$ suit une loi binomiale dont on précisera les paramètres.
  \item Calculer $P(X=3)$, puis $P(X\leq2)$ et $P(X\geq4)$.
  \item Calculer l'espérance, la variance et l'écart-type de $X$.
\end{questions}
\end{exercice}

\begin{exercice}[1]{Tests de logiciels}
Un développeur teste un algorithme sur $14$ jeux de données. Chaque jeu de données passe tous les tests avec la probabilité
$0{,}6$, indépendamment des autres. On note $X$ le nombre de jeux de données réussis.
\begin{questions}
  \item Donner la loi de $X$.
  \item Calculer $P(X=9)$ et $P(X\geq7)$.
  \item Calculer $E(X)$ et $\sigma(X)$.
  \item Chaque jeu de données réussi rapporte une prime de $100$\,€. Quel est le montant moyen de la prime ?
\end{questions}
\end{exercice}

\rubrique{Entraînement}

\begin{exercice}[2]{Un sondage}
Dans une ville de $5\,000$ habitants, 35\,\% des électeurs ont l'intention de voter pour un candidat. On interroge $25$ habitants
au hasard ; on assimile ce choix à un tirage avec remise. On note $X$ le nombre de réponses favorables.
\begin{questions}
  \item Justifier que $X$ suit une loi binomiale et donner ses paramètres.
  \item Calculer $P(X\geq10)$ et $P(X=8)$.
  \item Quel est le nombre moyen de réponses favorables ?
  \item Déterminer le plus petit entier $k$ tel que $P(X\leq k)\geq0{,}95$. Interpréter.
\end{questions}
\end{exercice}

\begin{exercice}[2]{Efficacité d'un traitement}
Un médicament contre la migraine est efficace dans 75\,\% des cas. Un médecin traite $12$ patients, dont les réactions sont
indépendantes. On note $X$ le nombre de patients guéris.
\begin{questions}
  \item Donner la loi de $X$, puis calculer $P(X=10)$ et $P(X\geq8)$.
  \item Quel est le nombre moyen de patients guéris ?
  \item Chaque traitement coûte $150$\,€. Quel est, en moyenne, le coût du traitement par patient guéri ?
\end{questions}
\end{exercice}

\begin{exercice}[2]{Génétique des pois}
On croise des pois hybrides. Chaque descendant a, indépendamment des autres, une chance sur deux d'avoir une gousse lisse. On
observe $16$ descendants et l'on note $X$ le nombre de gousses lisses.
\begin{questions}
  \item Donner la loi de $X$, puis calculer $P(X=9)$ et $P(X>10)$.
  \item Calculer $E(X)$, $V(X)$ et $\sigma(X)$, et interpréter $E(X)$.
\end{questions}
\end{exercice}

\begin{exercice}[2]{Satisfaction de la clientèle}
Une entreprise de vente en ligne a un taux de satisfaction de 82\,\%. Un commercial contacte $20$ clients au hasard ; on note $X$
le nombre de clients satisfaits.
\begin{questions}
  \item Donner la loi de $X$, puis calculer $P(X=16)$ et $P(X\geq18)$.
  \item Calculer l'écart-type de $X$.
  \item Chaque client satisfait rapporte $200$\,€ de chiffre d'affaires. Quel est le chiffre d'affaires moyen ?
\end{questions}
\end{exercice}

\rubrique{Problèmes}

\begin{exercice}[2]{Contrôle de qualité}
On contrôle $15$ produits ; chacun est défectueux avec la probabilité $0{,}1$, indépendamment des autres. On note $X$ le nombre
de produits défectueux.
\begin{questions}
  \item Calculer $P(X=2)$, $P(X\leq1)$ et $P(X\geq3)$.
  \item Calculer l'espérance et l'écart-type de $X$.
  \item Le remplacement d'un produit défectueux coûte $50$\,€. Quel est le coût moyen de remplacement ?
\end{questions}
\end{exercice}

\begin{exercice}[2]{Tests de résistance}
Un composant électronique présente un défaut avec la probabilité $0{,}08$. On teste un lot de $30$ composants, indépendants les uns
des autres, et l'on note $X$ le nombre de composants défectueux.
\begin{questions}
  \item Donner la loi de $X$, puis calculer $P(X\leq2)$ et $P(3\leq X\leq5)$.
  \item Calculer l'écart-type de $X$.
  \item Le remplacement d'un composant défectueux coûte $25$\,€. Quel est le coût moyen pour le lot ?
\end{questions}
\end{exercice}

\begin{exercice}[2]{Contrôle d'alcoolémie}
Lors d'un contrôle routier, 15\,\% des conducteurs dépassent le taux d'alcoolémie autorisé. On contrôle $40$ conducteurs choisis
au hasard ; on note $X$ le nombre de conducteurs en infraction.
\begin{questions}
  \item Donner la loi de $X$, puis calculer $P(X=6)$ et $P(X\leq4)$.
  \item Calculer la variance et l'écart-type de $X$.
  \item Chaque infraction coûte une amende de $135$\,€. Quel est le montant moyen des amendes ?
\end{questions}
\end{exercice}

\begin{exercice}[3]{Recyclage des déchets}
Dans une ville, 25\,\% des ménages trient correctement leurs déchets. On contrôle $35$ ménages choisis au hasard ; on note $X$ le
nombre de ménages qui trient correctement. Chacun d'eux fait économiser $200$\,€ à la collectivité.
\begin{questions}
  \item Donner la loi de $X$, puis calculer $P(X=10)$ et $P(6\leq X\leq12)$.
  \item Calculer la variance et l'écart-type de $X$.
  \item Une campagne de sensibilisation coûte $5\,000$\,€ et ferait passer le taux de tri de 25\,\% à 30\,\%. Calculer l'économie
    moyenne supplémentaire réalisée sur ces $35$ ménages. La campagne est-elle rentable à cette échelle ?
\end{questions}
\end{exercice}

\begin{exercice}[3]{Une campagne téléphonique}
Lors d'une campagne de démarchage, 15\,\% des personnes contactées acceptent un rendez-vous. Un employé passe $30$ appels par jour,
et l'on note $X$ le nombre de rendez-vous obtenus.
\begin{questions}
  \item Justifier que $X$ suit une loi binomiale dont on précisera les paramètres.
  \item Calculer la probabilité que l'employé obtienne exactement $5$ rendez-vous, puis au moins un rendez-vous.
  \item Calculer l'espérance et l'écart-type de $X$.
  \item Un autre employé contacte $n$ personnes. Déterminer le plus petit entier $n$ tel que la probabilité d'obtenir au moins
    un rendez-vous soit supérieure à $0{,}95$.
\end{questions}
\end{exercice}

\begin{exercice}[3]{Répondre au hasard à un QCM}
Un QCM comporte $20$ questions indépendantes ; chacune propose $4$ réponses, dont une seule est juste. Un élève répond au hasard à
toutes les questions ; on note $X$ le nombre de bonnes réponses.
\begin{questions}
  \item Donner la loi de $X$ et son espérance.
  \item Calculer $P(X=5)$ et $P(X\geq10)$. Que représente cette dernière probabilité ?
  \item Le barème attribue $1$ point par bonne réponse et retire $\frac13$ de point par mauvaise réponse. Exprimer la note $N$ en
    fonction de $X$, puis calculer $E(N)$. Pourquoi ce barème ?
  \item Calculer la probabilité que l'élève obtienne une note au moins égale à $5$.
\end{questions}
\end{exercice}

\begin{exercice}[3]{La surréservation des vols}
Un avion compte $200$ places. Chaque passager ayant acheté un billet se présente à l'embarquement avec la probabilité $0{,}95$,
indépendamment des autres. La compagnie vend $N$ billets, avec $N\geq200$ ; on note $X$ le nombre de passagers qui se présentent.
\begin{questions}
  \item On suppose $N=200$. Donner la loi de $X$ et son espérance, puis calculer $P(X=200)$. Interpréter.
  \item La compagnie vend $N=205$ billets. Quelle est la probabilité qu'au moins un passager ne puisse pas embarquer ?
  \item La compagnie accepte ce risque s'il ne dépasse pas 5\,\%. Que fait le programme ci-dessous ? Il affiche $205$ : interpréter.
\begin{python}
from math import comb

def risque(N):
    s = 0
    for k in range(201, N + 1):
        s = s + comb(N, k) * 0.95**k * 0.05**(N - k)
    return s

N = 200
while risque(N + 1) <= 0.05:
    N = N + 1
print(N)
\end{python}
\end{questions}
\end{exercice}

\begin{exercice}[3]{Une pièce est-elle équilibrée ?}
Pour savoir si une pièce est équilibrée, on la lance $100$ fois et l'on note $X$ le nombre de piles. Si la pièce est équilibrée,
$X$ suit la loi $\mathcal{B}(100\pv0{,}5)$. On donne : $P(X\leq39)\approx0{,}017\,6$, $P(X\leq40)\approx0{,}028\,4$,
$P(X\leq59)\approx0{,}971\,6$ et $P(X\leq60)\approx0{,}982\,4$.
\begin{questions}
  \item Calculer $E(X)$ et $\sigma(X)$.
  \item Déterminer le plus petit entier $a$ tel que $P(X\leq a)>0{,}025$ et le plus petit entier $b$ tel que $P(X\leq b)\geq0{,}975$.
    Calculer $P(a\leq X\leq b)$.
  \item On décide : « si $X\in\intFF ab$, la pièce est considérée comme équilibrée ; sinon, comme truquée ». Quelle est la
    probabilité de déclarer truquée une pièce équilibrée ?
  \item On obtient $65$ piles. Que conclure ? Et avec $58$ piles ?
  \item Avec $400$ lancers, la même méthode donne l'intervalle $\intFF{180}{220}$. Comparer les intervalles de fréquences
    obtenus pour $100$ et pour $400$ lancers. Commenter.
\end{questions}
\end{exercice}
